% GENERATED FILE. Edit canonical lessons, then run npm run generate:books.
% canonical-source-hash: fnv1a64:b4f9e108f3df7755
% canonical-lessons: GU-C325-ajhadi, GU-C325-gulam, GU-C325-angrej, GU-C325-adolan, GU-C325-balvo

\chapter{The Fight for Freedom}
\label{ch:gu-the-fight-for-freedom}
\hlchaptermodality{325}

\begin{chapteropening}
\textbf{By the end of this chapter:} \emph{I can say who fought for freedom and against whom.}

Complete the last lesson of chapter 325: I can say who fought for freedom and against whom.
\end{chapteropening}

\section[\texorpdfstring{\gu{આઝાદી} (ājhādī)}{આઝાદી (ājhādī)}]{\gu{આઝાદી} (ājhādī) — freedom}

\label{lesson:GU-C325-ajhadi}



\begin{quote}
Before the new one: say the Gujarati for control, then the Gujarati for to overthrow.
\end{quote}

\subsection*{You'll want to know: \gu{આઝાદી}}
\textbf{\gu{આઝાદી}} — \emph{ājhādī} — \textquotedblleft{}freedom\textquotedblright{}.

\begin{quote}
\gu{લોકો} \gu{આઝાદી} \gu{માટે} \gu{લડ્યા}.
\end{quote}

\begin{quote}
\emph{loko ājhādī māṭe laḍyā.}
\end{quote}

\begin{quote}
People fought for freedom.
\end{quote}

A feminine noun.

\begin{cousinweb}
\textbf{\gu{આઝાદી}} (\emph{ājhādī}) is usually traced to Persian \emph{āzādī}, \textquotedblleft{}freedom\textquotedblright{}, the noun of \emph{āzād}, free. It is feminine: \textbf{\gu{આઝાદી} \gu{માટે}}, for freedom.
\end{cousinweb}

\begin{grammarlens}[title={What you've built}]
One more word for telling what happened.
\end{grammarlens}

\subsection*{Your turn}
\begin{itemize}
\raggedright
  \item \emph{Say it:} \emph{ājhādī}
  \item \emph{Say it:} \emph{ājhādī}, once more
  \item \emph{Say it:} \emph{loko ājhādī māṭe laḍyā.} — People fought for freedom.
  \item \emph{Recall:} say the Gujarati for control, then the Gujarati for to overthrow, then say \emph{ājhādī} again
\end{itemize}

\begin{tcolorbox}[breakable,colback=teal!4,colframe=teal!35!black,title={Before you move on}]
What does \gu{આઝાદી} mean? (\textquotedblleft{}Freedom\textquotedblright{}.) Say it once more.
\end{tcolorbox}
\section[\texorpdfstring{\gu{ગુલામ} (gulām)}{ગુલામ (gulām)}]{\gu{ગુલામ} (gulām) — a slave}

\label{lesson:GU-C325-gulam}



\begin{quote}
Before the new one: say the Gujarati for to overthrow, then the Gujarati for freedom.
\end{quote}

\subsection*{You'll want to know: \gu{ગુલામ}}
\textbf{\gu{ગુલામ}} — \emph{gulām} — \textquotedblleft{}a slave\textquotedblright{}.

\begin{quote}
\gu{જૂના} \gu{વખતમાં} \gu{ઘણા} \gu{લોકો} \gu{ગુલામ} \gu{હતા}.
\end{quote}

\begin{quote}
\emph{jūnā vakhatmā̃ ghaṇā loko gulām hatā.}
\end{quote}

\begin{quote}
In the old days many people were slaves.
\end{quote}

A masculine or feminine noun, by the person meant.

\begin{cousinweb}
is usually traced, through Persian, to Arabic \emph{ġulām}, \textquotedblleft{}a boy, a servant\textquotedblright{}. In Gujarati it is a slave, and a people not free are spoken of as \textbf{\gu{ગુલામ}}.
\end{cousinweb}

\begin{grammarlens}[title={What you've built}]
One more word for telling what happened.
\end{grammarlens}

\subsection*{Your turn}
\begin{itemize}
\raggedright
  \item \emph{Say it:} \emph{gulām}
  \item \emph{Say it:} \emph{gulām}, once more
  \item \emph{Say it:} \emph{jūnā vakhatmā̃ ghaṇā loko gulām hatā.} — In the old days many people were slaves.
  \item \emph{Recall:} say the Gujarati for to overthrow, then the Gujarati for freedom, then say \emph{gulām} again
\end{itemize}

\begin{tcolorbox}[breakable,colback=teal!4,colframe=teal!35!black,title={Before you move on}]
What does \gu{ગુલામ} mean? (\textquotedblleft{}A slave\textquotedblright{}.) Say it once more.
\end{tcolorbox}
\section[\texorpdfstring{\gu{અંગ્રેજ} (angrej)}{અંગ્રેજ (angrej)}]{\gu{અંગ્રેજ} (angrej) — an Englishman; the English (rulers)}

\label{lesson:GU-C325-angrej}



\begin{quote}
Before the new one: say the Gujarati for freedom, then the Gujarati for a slave.
\end{quote}

\subsection*{You'll want to know: \gu{અંગ્રેજ}}
\textbf{\gu{અંગ્રેજ}} — \emph{angrej} — \textquotedblleft{}an Englishman; the English (rulers)\textquotedblright{}.

\begin{quote}
\gu{અંગ્રેજો} \gu{ઘણાં} \gu{વરસ} \gu{અહીં} \gu{રહ્યા}.
\end{quote}

\begin{quote}
\emph{angrejo ghaṇā̃ varas ahī̃ rahyā.}
\end{quote}

\begin{quote}
The English stayed here for many years.
\end{quote}

A masculine or feminine noun, by the person meant.

\begin{cousinweb}
\textbf{\gu{અંગ્રેજ}} (\emph{angrej}) is usually traced, through Persian \emph{angrez}, to Portuguese \emph{inglês}, \textquotedblleft{}English\textquotedblright{}. In the plural it is the English rulers of colonial times: \textbf{\gu{અંગ્રેજો}}.
\end{cousinweb}

\begin{grammarlens}[title={What you've built}]
One more word for telling what happened.
\end{grammarlens}

\subsection*{Your turn}
\begin{itemize}
\raggedright
  \item \emph{Say it:} \emph{angrej}
  \item \emph{Say it:} \emph{angrej}, once more
  \item \emph{Say it:} \emph{angrejo ghaṇā̃ varas ahī̃ rahyā.} — The English stayed here for many years.
  \item \emph{Recall:} say the Gujarati for freedom, then the Gujarati for a slave, then say \emph{angrej} again
\end{itemize}

\begin{tcolorbox}[breakable,colback=teal!4,colframe=teal!35!black,title={Before you move on}]
What does \gu{અંગ્રેજ} mean? (\textquotedblleft{}An Englishman; the English (rulers)\textquotedblright{}.) Say it once more.
\end{tcolorbox}
\section[\texorpdfstring{\gu{આંદોલન} (ā̃dolan)}{આંદોલન (ā̃dolan)}]{\gu{આંદોલન} (ā̃dolan) — a movement; a campaign}

\label{lesson:GU-C325-adolan}



\begin{quote}
Before the new one: say the Gujarati for a slave, then the Gujarati for an Englishman.
\end{quote}

\subsection*{You'll want to know: \gu{આંદોલન}}
\textbf{\gu{આંદોલન}} — \emph{ā̃dolan} — \textquotedblleft{}a movement; a campaign\textquotedblright{}.

\begin{quote}
\gu{ગામના} \gu{લોકોએ} \gu{મોટું} \gu{આંદોલન} \gu{કર્યું}.
\end{quote}

\begin{quote}
\emph{gāmnā lokoe moṭũ ā̃dolan karyũ.}
\end{quote}

\begin{quote}
The villagers started a big campaign.
\end{quote}

A neuter noun.

\begin{cousinweb}
\textbf{\gu{આંદોલન}} (\emph{ā̃dolan}) is a tatsama, Sanskrit \emph{āndolana}, \textquotedblleft{}swinging, a swing\textquotedblright{} (Monier-Williams); the sense of a stirring of people is the modern language's. It is neuter: \textbf{\gu{મોટું} \gu{આંદોલન}}.
\end{cousinweb}

\begin{grammarlens}[title={What you've built}]
One more word for telling what happened.
\end{grammarlens}

\subsection*{Your turn}
\begin{itemize}
\raggedright
  \item \emph{Say it:} \emph{ā̃dolan}
  \item \emph{Say it:} \emph{ā̃dolan}, once more
  \item \emph{Say it:} \emph{gāmnā lokoe moṭũ ā̃dolan karyũ.} — The villagers started a big campaign.
  \item \emph{Recall:} say the Gujarati for a slave, then the Gujarati for an Englishman, then say \emph{ā̃dolan} again
\end{itemize}

\begin{tcolorbox}[breakable,colback=teal!4,colframe=teal!35!black,title={Before you move on}]
What does \gu{આંદોલન} mean? (\textquotedblleft{}A movement; a campaign\textquotedblright{}.) Say it once more.
\end{tcolorbox}
\section[\texorpdfstring{\gu{બળવો} (baḷvo)}{બળવો (baḷvo)}]{\gu{બળવો} (baḷvo) — a revolt}

\label{lesson:GU-C325-balvo}



\begin{quote}
Before the new one: say the Gujarati for an Englishman, then the Gujarati for a campaign.
\end{quote}

\subsection*{You'll want to know: \gu{બળવો}}
\textbf{\gu{બળવો}} — \emph{baḷvo} — \textquotedblleft{}a revolt\textquotedblright{}.

\begin{quote}
\gu{એક} \gu{રાત્રે} \gu{સૈનિકોનો} \gu{બળવો} \gu{શરૂ} \gu{થયો}.
\end{quote}

\begin{quote}
\emph{ek rātre sainikono baḷvo śarū thayo.}
\end{quote}

\begin{quote}
One night the soldiers' revolt began.
\end{quote}

A masculine noun.

\begin{cousinweb}
\textbf{\gu{બળવો}} (\emph{baḷvo}) No dictionary we use records an older form, so its history is not settled. It is masculine: \textbf{\gu{બળવો} \gu{થયો}}, there was a revolt; \textbf{\gu{બળવો} \gu{કર્યો}}, they rose up.
\end{cousinweb}

\begin{grammarlens}[title={What you've built}]
One more word for telling what happened.
\end{grammarlens}

\subsection*{Your turn}
\begin{itemize}
\raggedright
  \item \emph{Say it:} \emph{baḷvo}
  \item \emph{Say it:} \emph{baḷvo}, once more
  \item \emph{Say it:} \emph{ek rātre sainikono baḷvo śarū thayo.} — One night the soldiers' revolt began.
  \item \emph{Recall:} say the Gujarati for an Englishman, then the Gujarati for a campaign, then say \emph{baḷvo} again
\end{itemize}

\begin{tcolorbox}[breakable,colback=teal!4,colframe=teal!35!black,title={Before you move on}]
What does \gu{બળવો} mean? (\textquotedblleft{}A revolt\textquotedblright{}.) Say it once more.
\end{tcolorbox}
